%%%%%%%%%%%%
%
% $Autor: Wings $
% $Datum: 2019-03-05 08:03:15Z $
% $Pfad: TemplateSensor $
% $Version: 4250 $
% !TeX spellcheck = en_GB/de_DE
% !TeX encoding = utf8
% !TeX root = filename 
% !TeX TXS-program:bibliography = txs:///biber
%
%%%%%%%%%%%%

\chapter{LoRaWAN-Architektur und Datenfluss}
	Dieses Kapitel beschreibt den grundlegenden Aufbau der verwendeten \acs{lorawan}-Architektur sowie den Datenfluss eines Uplinks vom Endgerät über das Gateway bis zum Netzwerkserver.

\section{LoRaWAN-Architektur}
	Die \acs{lorawan}-Architektur besteht aus drei zentralen Komponenten: den Endgeräten, den Gateways sowie dem \acl{lns} (siehe~\autoref{fig:lorawanArchitectureGeneral}). Diese Komponenten sind funktional klar voneinander getrennt und übernehmen jeweils spezifische Aufgaben innerhalb des Gesamtsystems.

	\begin{figure}[H]
		\centering
		\begin{tikzpicture}
			\tikzset{
				node distance=0.5cm and 0.4cm,
				textDef/.style={font=\sffamily\small,},
				lineDef/.style={line width={2pt}},
				arrowStyle/.style={line width=2pt, >=stealth, red},
			}
			\sensorEndNode{1}{0}{-7}{0.2}{E1}
			\gatewayNode{1}{0}{-3}{1.45}{G1}
			\serverNode{0.6}{0}{0}{0}{S1}
			
			\draw[arrowStyle, <->] (E1.east) -- (G1.west) node[textDef, midway, above] {\shortstack{LoRa\\Modulation}};
			\draw[arrowStyle, <->] (G1.east) -- (S1.west) node[textDef, midway, above] {\shortstack{\acs{wan}\\Backhaul}};
			
			\node[textDef, below= of E1.south] {Endgerät};
			\node[textDef, below= of G1.south] {Gateway};
			\node[textDef, below= of S1.south] {\shortstack{LoRaWAN\\Server}};
			
		\end{tikzpicture}
		\captionof{figure}{Kommunikationspfad eines Uplinks im LoRaWAN}
		\label{fig:lorawanArchitectureGeneral}
	\end{figure}
	
	Die Kommunikation zwischen Endgerät und Gateway erfolgt drahtlos unter Verwendung der LoRa-Modulation im Sub-GHz-Frequenzbereich (siehe~\autoref{lora:general}). Zwischen Gateway und \acs{lns} findet die Datenübertragung über eine IP-basierte \acs{wan}-Verbindung statt, beispielsweise über Ethernet oder Mobilfunk (siehe~\autoref{gateway:backhaul}). Das Gateway fungiert dabei ausschlie{\ss}lich als Weiterleitungsinstanz und führt keine inhaltliche Verarbeitung der Nutzdaten durch.

\subsection{Ablauf eines Uplinks}
	Nach der Erfassung der Sensordaten erzeugt das Endgerät einen \ac{lorawan}-Uplink-Frame, welcher verschlüsselt und über die LoRa-Funkschnittstelle ausgesendet wird. Ein oder mehrere Gateways können diesen Frame empfangen und leiten ihn unverändert an den zuständigen \acs{lns} weiter (siehe~\autoref{gateway:info}).
	
	Das Ablaufdiagramm in Abbildung~\ref{LoRaWAN:Architecture:uplink} zeigt den vollständigen Weg eines Uplinks vom Endgerät bis zum \acs{lns}. Dabei werden die vom Sensor erfassten Benutzerdaten im Feld \PYTHON{FRMPayload} des \ac{lorawan}-Frames übertragen (siehe~\autoref{lorawan:frmpayload}).
	
	\begin{figure}[H]
		\centering
		\includegraphics[width=0.8\textwidth]{LoRaAndLoRaWAN/Flowcharts/SensorsAPI/SensorsAPIUplinkProcess}
		\caption{Ablaufdiagramm eines Uplinks vom Endgerät bis zum \ac{lns}}
		\label{LoRaWAN:Architecture:uplink}
	\end{figure}
	


\subsection{Ablauf der Bearbeitung eines Uplinks auf dem ChirpStack Network Server}
	Der ChirpStack Network Server ist für die vollständige Verarbeitung aller eingehenden \ac{lorawan}-Frames verantwortlich. Dazu zählen unter anderem die Authentifizierung des Endgeräts, die Entschlüsselung der Nutzdaten, die Verwaltung der Frame Counter sowie die Duplikaterkennung (siehe~\autoref{LoRaWAN:Architecture:LNS}).
	
	Abbildung~\ref{LoRaWAN:Architecture:LNS} stellt den Ablauf der Verarbeitung eines Uplinks mit Applikationsdaten auf dem ChirpStack Network Server dar.

	
	\begin{figure}[H]
		\centering
		\includegraphics[width=0.8\textwidth]{LoRaAndLoRaWAN/Flowcharts/SensorsAPI/SensorsAPIUplinkProcessOnServer.pdf}
		\caption{Ablaufdiagramm der Verarbeitung eines Uplinks im \ac{lns}}
		\label{LoRaWAN:Architecture:LNS}
	\end{figure}
	
	Nach erfolgreicher Verarbeitung werden die entschlüsselten Nutzdaten an die Anwendung des Benutzers weitergeleitet, wo sie beispielsweise gespeichert, visualisiert oder weiterverarbeitet werden können (siehe~\autoref{server:appserver}).
